\tikzsetfigurename{14_arch-}
% Part V packs many tables and figures: allow fuller float areas (affects only Part V and the appendices).
\setcounter{topnumber}{3}\setcounter{bottomnumber}{2}\setcounter{totalnumber}{5}
\renewcommand{\topfraction}{0.9}\renewcommand{\bottomfraction}{0.6}
\renewcommand{\textfraction}{0.08}\renewcommand{\floatpagefraction}{0.8}

\chapter{The Open-SERDES Architecture}\label{ch:arch}

\chapterlead{Every decision of the last thirteen chapters on a few pages: the blocks, the clocks, the number each block must meet, and what it all costs in power.}

\begin{objectives}
\begin{itemize}
  \item Read the complete \OS{} block diagram and follow a bit from the 16-bit TX word to the 16-bit RX word.
  \item Trace every clock in the chip from the 50\,MHz reference to the last divider, with its frequency and phase.
  \item Use the block specification table to find the target, and the chapter that derives it, for any block you design.
  \item Build a first-cut power budget from tail currents and convert it to energy per bit.
  \item Defend the ten main architecture decisions and name the risks that remain open.
\end{itemize}
\end{objectives}

%% =====================================================================
\section{The complete transceiver}\label{sec:arch-block}

\Cref{fig:arch} puts the whole chip on one page. It is a single-lane transceiver: one transmitter, one receiver and one clock generator on the same \sgx{} die, connected to each other through the off-chip REF channel (\cref{sec:refchannel}) or through an on-chip loopback.

\paragraph{Transmit path} A 16-bit word arrives every 3.2\,ns (312.5\,MHz), from the user or from the on-chip PRBS7/31 generator (\cref{sec:prbs}). CMOS on 1.2\,V performs the first 16:4 serialization; HBT CML on the 2.5\,V analog supply performs the 4:2 and the final half-rate 2:1, built three times so the driver receives $D_{n+1}$, $D_n$ and $D_{n-1}$ for its three FFE taps (\cref{sec:serializer,sec:ffe-impl}). The CML driver launches 1\,Vppd into \qty{50}{\ohm} per side (\cref{sec:cml-driver}).

\paragraph{Receive path} Behind the ESD diodes and the \qty{100}{\ohm} termination (\cref{sec:esd-pad}), a loopback multiplexer selects the pad or the near-end loopback. The CTLE adds 0--12\,dB of peaking (\cref{ch:ctle}), the DFE summer subtracts $w_1$ and $w_2$ times the previous decisions (\cref{sec:dfe-timing}), and four samplers decide: data D0, D1 and edges E0, E1 (\cref{sec:sampler}). The deserializer delivers 16 data and 16 edge bits at 312.5\,MHz (\cref{sec:deserializer}); the PRBS checker counts errors.

\paragraph{Clocking and control} An integer-N PLL with an LC VCO multiplies the 50\,MHz reference by $N=50$ to 2.5\,GHz (\cref{ch:pll}). The TX gets it through buffers and $\div2$ dividers; the RX gets I and Q from a polyphase filter (\cref{sec:quadrature}), rotated by two phase interpolators (PIs) under the digital CDR loop filter (\cref{sec:picdr}). An SPI port writes every setting.

\begin{figure}[!tb]
\centering
% ch14_arch.tex -- fig:arch : the complete Open-SERDES transceiver (Chapter 14)
\begin{tikzpicture}[x=1cm,y=1cm,font=\sffamily\scriptsize,
  a/.style={blkb,minimum height=9mm,minimum width=13mm,font=\sffamily\scriptsize,inner sep=1.5pt},
  d/.style={blkg,minimum height=9mm,minimum width=13mm,font=\sffamily\scriptsize,inner sep=1.5pt},
  rt/.style={font=\sffamily\scriptsize,text=muted,align=center},
  rate/.style={font=\sffamily\scriptsize\bfseries,text=slate,align=center},
  ck/.style={font=\sffamily\scriptsize,text=golddk,align=center},
  cfg/.style={-{Stealth[length=1.6mm]},draw=muted,line width=0.6pt,dotted},
  lb/.style={-{Stealth[length=2.2mm]},draw=c5,line width=0.9pt},
  pad/.style={draw=navy,fill=navy,minimum size=2.2mm,inner sep=0pt},
  bandlbl/.style={font=\sffamily\scriptsize\bfseries,text=slate,anchor=north west}]
% ------------------------------------------------------------ die and bands
\draw[navy,line width=1.1pt,rounded corners=5pt,fill=offwhite] (0,0.05) rectangle (13.0,9.1);
\fill[gtint!35] (0.08,7.22) rectangle (12.92,9.02);
\fill[bgpale]   (0.08,2.82) rectangle (12.92,7.14);
\fill[btint!40] (0.08,0.13) rectangle (12.92,2.74);
\node[bandlbl] at (0.12,9.0) {TRANSMITTER};
\node[bandlbl,anchor=south west] at (0.12,2.9) {CLOCKING AND CDR};
\node[bandlbl] at (0.12,2.72) {RECEIVER};
\node[font=\sffamily\scriptsize\bfseries,text=navy,anchor=south east] at (13.0,9.12) {\OS{} die, IHP \sgx{}};
% ------------------------------------------------------------ TX row (y = 8.0)
\node[d,minimum width=18mm] (prbs) at (1.45,8.0) {PRBS7/31 gen.\\+ data mux};
\node[d,minimum width=13mm] (s164) at (3.95,8.0) {16:4\\CMOS};
\node[a,minimum width=12mm] (s42)  at (5.95,8.0) {4:2\\CML};
\node[a,minimum width=16mm] (s21)  at (8.0,8.0) {2:1 CML $\times$3\\half-rate};
\node[a,minimum width=19mm,line width=1pt] (drv) at (10.55,8.0) {\textbf{3-tap FFE}\\CML driver};
\node[pad] (txpad) at (13.0,8.0) {};
\draw[sig,line width=1.1pt] (-0.8,8.0) -- (prbs.west);
\node[rt,anchor=south] at (-0.4,8.05) {16\,b\\in};
\draw[sig] (prbs) -- node[rt,above=0pt]{16} (s164);
\draw[sig] (s164) -- node[rt,above=0pt]{4} (s42);
\draw[sig] (s42) -- node[rt,above=0pt]{2} (s21);
\draw[sig] (s21) -- node[rt,above=0pt]{3} (drv);
\foreach \x/\r in {2.82/{312.5\,Mb/s}, 4.97/{1.25\,Gb/s}, 6.95/{2.5\,Gb/s}, 9.2/{5\,Gb/s}}{\node[rate] at (\x,8.72) {\r};}
\draw[navy,line width=1.1pt] (drv.east) -- (txpad.west); \coordinate (lbtap) at (12.25,8.0);
\node[rt,anchor=south] at (12.6,8.05) {TX$\pm$};
% ------------------------------------------------------------ channel (off-chip)
\node[blk,minimum width=16mm,minimum height=46mm,fill=white,font=\sffamily\scriptsize,align=center] (chan) at (14.15,4.8)
  {\textbf{Channel}\\[3pt]REF: 1\,m FR-4\\+ 300\,fF pads\\[3pt]$-17.4$\,dB\\at 2.5\,GHz\\[3pt]\textcolor{muted}{off-chip:}\\\textcolor{muted}{SMA, PCB}\\\textcolor{muted}{or cable}};
\draw[sig,line width=1.1pt] (txpad.east) -| (chan.north);
\draw[sig,line width=1.1pt] (chan.south) |- (13.0,1.6);
% ------------------------------------------------------------ RX row (y = 1.6), signal flows right to left
\node[pad] (rxpad) at (13.0,1.6) {};
\node[a,minimum width=14mm] (term) at (11.8,1.6) {\qty{100}{\ohm} term.\\+ ESD};
\node[a,minimum width=8mm,minimum height=7mm] (lbm) at (10.5,1.6) {LB\\mux};
\node[a,minimum width=13mm] (ctle) at (9.3,1.6) {CTLE\\0--12\,dB};
\node[sumnode,minimum size=6.5mm] (sum) at (8.05,1.6) {$\Sigma$};
\node[a,minimum width=16.5mm] (samp) at (6.55,1.6) {4 samplers\\D0 E0 D1 E1};
\node[d,minimum width=13mm] (des) at (4.45,1.6) {1:16\\deser.};
\node[d,minimum width=14mm] (chk) at (2.3,1.6) {PRBS\\checker};
\node[rt,anchor=south] at (12.76,1.63) {RX$\pm$};
\draw[navy,line width=1.1pt] (rxpad.west) -- (term.east);
\draw[sig] (term) -- (lbm);
\draw[sig] (lbm) -- (ctle);
\draw[sig] (ctle) -- (sum);
\draw[sig] (sum) -- (samp);
\node[rt,anchor=south,text=bblue] at (8.05,1.95) {DFE};
\draw[sig] (samp) -- node[rt,above=0pt]{4} (des);
\node[rate,anchor=north] at (5.42,1.1) {2.5\,Gb/s};
\draw[sig] (des) -- node[rt,above=0pt]{32} (chk);
\node[rate,anchor=north] at (3.37,1.1) {312.5\,M};
\draw[sig,line width=1.1pt] (chk.west) -- (-0.8,1.6);
\node[rt,anchor=south] at (-0.4,1.65) {16\,b\\out};
% DFE feedback: decisions -> taps w1, w2 -> summer
\draw[sigb] ([xshift=4mm]samp.south) -- ++(0,-0.52) -| (sum.south);
\node[font=\sffamily\scriptsize,text=bblue,anchor=north] at (7.6,0.6) {taps $w_1$, $w_2$};
% ------------------------------------------------------------ clocking row (y = 5.1)
\node[a,minimum width=19mm] (pll) at (1.6,5.1) {\textbf{PLL}, $N=50$\\LC VCO 2.5\,GHz};
\node[a,minimum width=17mm] (ppf) at (4.25,5.1) {I/Q polyphase\\+ CML buffers};
\node[a,minimum width=21mm] (pi)  at (6.7,5.1) {2 phase interp.\\64 steps (6.25\,ps)};
\node[a,minimum width=55mm,minimum height=5.5mm] (txck) at (5.9,6.55) {TX clock buffers and $\div2$, $\div4$, $\div8$ dividers};
\node[d,minimum width=18mm,minimum height=7.5mm] (lf) at (4.45,3.6) {CDR loop filter\\P + I, 625\,MHz};
\node[d,minimum width=19mm] (spi) at (10.4,5.1) {SPI +\\config registers};
% refclk
\draw[clk] (-0.8,5.1) -- (pll.west);
\node[pad] at (0,5.1) {};
\node[ck,anchor=south] at (-0.4,5.15) {REF\\50\,M};
% PLL -> TX clock tree -> serializer
\draw[clk] (pll.north) |- node[ck,pos=0.3,left=0pt]{2.5\,GHz} (txck.west);
\draw[clk] (3.35,6.55 |- txck.north) -- (3.35,7.05) -| (prbs.south);
\node[ck,anchor=east] at (1.4,7.3) {312.5\,MHz};
\draw[clk] (txck.north -| s164.south) -- node[ck,right=0pt,pos=0.45]{625\,M} (s164.south);
\draw[clk] (txck.north -| s42.south) -- node[ck,right=0pt,pos=0.45]{1.25\,G} (s42.south);
\draw[clk] (txck.north -| s21.south) -- node[ck,right=0pt,pos=0.45]{2.5\,G} (s21.south);
% PLL -> PPF -> PI -> samplers
\draw[clk] (pll.east) -- (ppf.west);
\draw[clk] (ppf.east) -- node[ck,above=0pt]{I, Q} (pi.west);
\draw[clk] (7.15,5.1 |- pi.south) -- coordinate[pos=0.78] (rck) node[ck,pos=0.3,right=0pt]{\ang{0}/\ang{90}/\ang{180}/\ang{270}} (7.15,1.6 |- samp.north);
% recovered clock to the deserializer
\fill[gold] (rck) circle (1.3pt);
\draw[clk] (rck) -| node[ck,pos=0.25,above=0pt]{recovered $\div2^n$} (4.85,1.6 |- des.north);
% CDR: deserializer -> loop filter -> PI code
\draw[sigg] (4.05,1.6 |- des.north) -- node[rt,left=0pt]{edge + data} (4.05,3.6 |- lf.south);
\draw[sigg] (lf.east) -| node[rt,pos=0.25,above=0pt]{code $k$} (6.2,5.1 |- pi.south);
% configuration (dotted)
\draw[cfg] (spi.north) -- node[rt,right=0pt,pos=0.5]{FFE taps} (spi.north |- drv.south);
\draw[cfg] (9.7,5.1 |- spi.south) -- node[rt,left=0pt,pos=0.55]{peaking} (9.7,1.6 |- ctle.north);
\draw[cfg] (spi.west) -- node[rt,above=0pt]{override} (pi.east);
\node[rt,anchor=north west,align=left,text width=19mm] at (10.0,4.45) {also DFE and offset DACs, PLL band, PRBS mode, loopback};
% near-end loopback
\fill[c5] (lbtap) circle (1.5pt);
\draw[lb] (lbtap) -- (12.25,2.4) -| (lbm.north);
\node[font=\sffamily\scriptsize,text=c5,rotate=90,anchor=south] at (12.25,6.1) {near-end loopback};
% ------------------------------------------------------------ legend
\begin{scope}[shift={(0.3,-0.45)}]
  \node[blkb,minimum width=6mm,minimum height=3.2mm,inner sep=0pt] at (0,0) {};
  \node[rt,anchor=west,text=navy] at (0.4,0) {analog / HBT CML (2.5\,V)};
  \node[blkg,minimum width=6mm,minimum height=3.2mm,inner sep=0pt] at (3.85,0) {};
  \node[rt,anchor=west,text=navy] at (4.25,0) {digital CMOS (1.2\,V)};
  \draw[clk] (7.25,0) -- (8.05,0); \node[rt,anchor=west,text=navy] at (8.1,0) {clock};
  \draw[lb] (9.05,0) -- (9.85,0); \node[rt,anchor=west,text=navy] at (9.9,0) {loopback};
  \draw[cfg] (11.25,0) -- (12.05,0); \node[rt,anchor=west,text=navy] at (12.1,0) {configuration};
\end{scope}
\end{tikzpicture}

\caption{\textbf{The complete \OS{} transceiver: a CMOS core at 312.5\,MHz, wrapped in HBT CML wherever a signal or clock runs at 1.25\,GHz or faster.} Read the TX band left to right and the RX band right to left; bold numbers above the TX row are lane rates, and slash numbers count parallel lanes. Notice three things. The 2.5\,GHz PLL clock is used twice: divided down for the TX (fixed phase), and split into I/Q and rotated by the PIs for the RX (recovered phase). The CDR is a digital loop: edge and data bits leave the deserializer, the loop filter computes the PI code, and the samplers move. The near-end loopback (violet) takes the driver output straight to the RX input, so the whole link can be tested without a board channel.}
\label{fig:arch}
\end{figure}

\subsection{The clocking plan}\label{sec:ch14-clock}

\Cref{fig:ch14-clock} follows each clock from the reference to the last flip-flop. All of the frequencies come from one relation between the bit rate, the half-rate clock and the word rate.

\begin{keyeq}[The \OS{} clock frequencies]
\begin{equation}
  f_\mathrm{clk} = \frac{\fb}{2} = N f_\mathrm{ref} = \qty{2.5}{\giga\hertz},\qquad
  f_\mathrm{word} = \frac{\fb}{16} = \qty{312.5}{\mega\hertz},\qquad
  \Delta t_\mathrm{PI} = \frac{1}{64\,f_\mathrm{clk}} = \qty{6.25}{\pico\second}
  \label{eq:ch14-clk}
\end{equation}
\end{keyeq}

The TX needs one clock phase: the final 2:1 mux uses both edges of the 2.5\,GHz clock as its select (\cref{sec:serializer}). The RX needs four phases half a UI apart: data samplers on the edges of I (eye centers), edge samplers on the edges of Q (crossings), all moved together by the CDR in 6.25\,ps steps. The RX dividers run from the \emph{recovered} clock, so the deserializer and the CDR logic sit in a clock domain that is not phase-locked to the TX word clock; a small FIFO crosses back to the core (\cref{sec:deserializer}). In near-end loopback both ends share the PLL, so the frequency offset is zero. Through a board to a second chip, the CDR must absorb up to $\pm300$\,ppm (\cref{sec:picdr}).

\begin{figure}[!tbp]
\centering
% ch14_clock.tex -- fig:ch14-clock : Open-SERDES clocking plan (frequencies and phases)
\begin{tikzpicture}[x=0.94cm,y=1cm,font=\sffamily\scriptsize,
  a/.style={blkb,minimum height=7.5mm,minimum width=12mm,font=\sffamily\scriptsize,inner sep=1.5pt},
  dv/.style={blkb,minimum height=6mm,minimum width=8mm,font=\sffamily\scriptsize,inner sep=1pt},
  fq/.style={font=\sffamily\scriptsize\bfseries,text=golddk,align=center},
  ds/.style={font=\sffamily\scriptsize,text=navy,align=center},
  tp/.style={circle,fill=gold,inner sep=0pt,minimum size=3pt}]
% ---------------- PLL core
\node[fq,anchor=east] (ref) at (0.35,0) {REFCLK\\50\,MHz\\(20\,ns)};
\node[a,minimum width=16mm] (pfd) at (1.4,0) {PFD + CP\\+ loop filter};
\node[a,minimum width=14mm] (vco) at (3.8,0) {LC VCO\\4\,nH, $Q$ = 16.8};
\node[dv] (div50) at (2.6,-1.05) {$\div50$};
\draw[clk] (ref.east) -- (pfd.west);
\draw[sig,draw=navy] (pfd.east) -- node[ds,above=0pt]{$V_\mathrm{ctrl}$} (vco.west);
\coordinate (v) at (5.0,0);
\draw[clk] (vco.east) -- (v);
\fill[gold] (v) circle (1.6pt);
\draw[clk] (v) |- (div50.east);
\draw[clk] (div50.west) -| (pfd.south);
\node[fq,anchor=west,align=left] at (5.1,0) {2.5\,GHz\\\textcolor{muted}{400\,ps}};
% ---------------- TX branch (y = 1.45)
\node[a] (tbuf) at (6.2,1.45) {CML\\buffer};
\node[dv] (t1) at (8.55,1.45) {$\div2$};
\node[dv] (t2) at (10.85,1.45) {$\div2$};
\node[dv] (t3) at (13.15,1.45) {$\div2$};
\draw[clk] (v) |- (tbuf.west);
\draw[clk] (tbuf) -- coordinate[pos=0.3] (tp0) node[fq,pos=0.66,above=1pt]{2.5\,GHz} (t1);
\draw[clk] (t1) -- coordinate[pos=0.3] (tp1) node[fq,pos=0.64,above=1pt]{1.25\,GHz} (t2);
\draw[clk] (t2) -- coordinate[pos=0.3] (tp2) node[fq,pos=0.64,above=1pt]{625\,MHz} (t3);
\draw[clk] (t3) -- coordinate[pos=0.3] (tp3) node[fq,pos=0.68,above=1pt]{312.5\,MHz} (15.45,1.45);
\foreach \p/\t in {tp0/{final 2:1 select\\(both edges: 5\,Gb/s)}, tp1/{4:2 CML,\\CMOS$\to$CML}, tp2/{CMOS 8:4}, tp3/{CMOS 16:8,\\TX word clock}}{
  \node[tp] at (\p) {}; \draw[clk,-] (\p) -- ++(0,-0.35); \node[ds,anchor=north] at ($(\p)+(0,-0.35)$) {\t};}
\node[lblb,anchor=west] at (5.55,2.3) {TX: one fixed phase, divided down};
% ---------------- RX branch (y = -1.75)
\node[a,minimum width=15mm] (ppf) at (6.1,-1.75) {2-stage PPF\\+ CML buffers};
\node[a,minimum width=16mm] (pi) at (8.6,-1.75) {PI$_D$: code $k$\\PI$_E$: $k+16$};
\node[dv] (r1) at (10.75,-1.75) {$\div2$};
\node[dv] (r2) at (12.4,-1.75) {$\div2$};
\node[dv] (r3) at (14.05,-1.75) {$\div2$};
\draw[clk] (v) |- (ppf.west);
\draw[clk] (ppf) -- coordinate (mq) (pi);
\node[ds,text=golddk,anchor=north] at ($(mq)+(0,-0.45)$) {I, Q, $\bar{\mathrm{I}}$, $\bar{\mathrm{Q}}$};
\draw[clk] (pi) -- coordinate[pos=0.25] (rp0) (r1);
\draw[clk] (r1) -- coordinate[pos=0.25] (rp1) (r2);
\draw[clk] (r2) -- coordinate[pos=0.25] (rp2) (r3);
\draw[clk] (r3) -- coordinate[pos=0.25] (rp3) (15.45,-1.75);
\foreach \p/\t/\dx in {rp0/{\textcolor{golddk}{\bfseries 2.5\,GHz}\\samplers D0 D1\\and E0 E1}/-0.1, rp1/{\textcolor{golddk}{\bfseries 1.25\,GHz}\\CMOS 1:2}/0, rp2/{\textcolor{golddk}{\bfseries 625\,MHz}\\CDR loop\\(8\,UI)}/0, rp3/{\textcolor{golddk}{\bfseries 312.5\,MHz}\\1:16 + FIFO}/0.25}{
  \node[tp] at (\p) {}; \draw[clk,-] (\p) -- ++(0,-0.35); \node[ds,anchor=north] at ($(\p)+(\dx,-0.35)$) {\t};}
\node[lblb,anchor=west] at (5.4,-0.85) {RX: four phases, rotated by the CDR};
% ---------------- timing of the four sampling phases (lower panel)
\begin{scope}[shift={(4.9,-5.05)},x=0.0185cm]
  \foreach \c/\n in {0/{bit $n$ (even)},200/{bit $n{+}1$ (odd)},400/{bit $n{+}2$}}{
    \draw[navy,line width=0.8pt,fill=btint] (\c-100,0.55) -- (\c-75,0.95) -- (\c+75,0.95) -- (\c+100,0.55) -- (\c+75,0.15) -- (\c-75,0.15) -- cycle;
    \node[ds] at (\c,0.55) {\n};}
  \draw[navy,line width=0.8pt] (-150,0.95) -- (-125,0.95) -- (-100,0.55) -- (-125,0.15) -- (-150,0.15);
  \node[ds,anchor=east] at (-160,0.55) {data at sampler};
  \draw[golddk,line width=0.9pt] (-150,-0.55) -- (0,-0.55) -- (0,-0.05) -- (200,-0.05) -- (200,-0.55) -- (400,-0.55) -- (400,-0.05) -- (550,-0.05);
  \node[ds,anchor=east] at (-160,-0.3) {clock I (\ang{0})};
  \draw[golddk,line width=0.9pt,dashed] (-150,-1.35) -- (100,-1.35) -- (100,-0.85) -- (300,-0.85) -- (300,-1.35) -- (550,-1.35);
  \node[ds,anchor=east] at (-160,-1.1) {clock Q (\ang{90})};
  \foreach \t/\s in {0/D0,100/E0,200/D1,300/E1,400/D0}{
    \draw[c4,densely dotted,line width=0.7pt] (\t,1.1) -- (\t,-1.55);
    \node[font=\sffamily\scriptsize\bfseries,text=c4,anchor=north] at (\t,-1.55) {\s};}
  \draw[<->,muted,line width=0.5pt] (0,1.25) -- node[ds,above=0pt]{1\,UI = 200\,ps} (200,1.25);
  \draw[<->,muted,line width=0.5pt] (0,-2.05) -- node[ds,below=0pt]{$T_\mathrm{clk}$ = 400\,ps = 2\,UI = 64 PI steps of 6.25\,ps} (400,-2.05);
\end{scope}
\end{tikzpicture}

\caption{\textbf{One 2.5\,GHz PLL serves the whole chip; the TX divides it, the RX rotates it.} Top: the frequency at every node. The TX chain divides by 2, 4 and 8 down to the 312.5\,MHz word clock. The RX chain first makes four phases with the polyphase filter, rotates them with the two PIs, and divides the recovered clock for the deserializer and for the CDR loop, which updates every 8\,UI at 625\,MHz. Bottom: what the four phases do. D0 and D1 sample the eye centers on the two edges of I; E0 and E1 sample the crossings on the edges of Q. One 400\,ps clock period carries two bits and contains 64 PI steps.}
\label{fig:ch14-clock}
\end{figure}

\begin{insight}
Three boundaries organize the chip, and most integration bugs live on them: the \textbf{speed boundary} between CMOS and CML (1.25\,Gb/s in the TX, 2.5\,Gb/s in the RX), which needs level converters and half a slow-clock cycle of timing margin; the \textbf{supply boundary} between 1.2 and 2.5\,V, where every $V_\mathrm{CE}$ must stay below $\mathrm{BV}_\mathrm{CEO}$; and the \textbf{clock-domain boundary} between the TX clock and the recovered RX clock. Draw all three on the floorplan first.
\end{insight}

%% =====================================================================
\section{Block specifications}\label{sec:ch14-specs}

\Cref{tab:specs} is the contract between the system model and the circuit designers: one row per block, the numbers that block must meet, and the chapter where each number is derived. Values marked $^\dagger$ are first estimates from hand analysis or behavioral models. They must be confirmed by simulation in \sgx{}, and some of them will move.

\begin{table}[!tbp]
\centering
\caption{\textbf{\OS{} block specifications.} $^\dagger$\,First estimate, not yet confirmed by circuit simulation in \sgx{}.}
\label{tab:specs}
\footnotesize
\renewcommand{\arraystretch}{1.12}
\begin{tabularx}{\linewidth}{@{}>{\raggedright\arraybackslash}p{19mm} >{\raggedright\arraybackslash}p{27mm} L >{\raggedright\arraybackslash}p{19mm}@{}}
\toprule
\textbf{Block} & \textbf{Key specification} & \textbf{Target value} & \textbf{Derived in}\\
\midrule
Link & rate, BER, channel & 5\,Gb/s NRZ; BER $\le10^{-12}$ over REF ($-17.4$\,dB at 2.5\,GHz); margin up to 20\,dB & \cref{sec:os-targets,sec:refchannel}\\
Serializer & ratio, clocks, jitter & 16:1 from 16\,b at 312.5\,MHz; CMOS 16:4, CML 4:2, half-rate CML 2:1; DCD + DDJ $\le8$\,ps$^\dagger$ & \cref{sec:serializer}\\
TX driver & topology, swing & HBT CML, \qty{50}{\ohm} back-termination per side; 0.8--1.0\,Vppd ($I_T$ = 16--20\,mA); every $V_\mathrm{CE}<\mathrm{BV}_\mathrm{CEO}\approx1.6$\,V at all corners & \cref{sec:cml-driver}\\
TX FFE & taps, range, resolution & $c_{-1}, c_0, c_1$ with $\sum\lvert c_k\rvert=1$; $c_1$ to $-0.25$ (6\,dB), $c_{-1}$ to $-0.1$; 5-bit current cells (steps of 1/32) & \cref{sec:ffe-taps,sec:ffe-impl}\\
Termination, pad & impedance, $C_p$, return loss & \qty{100}{\ohm} diff., $\pm10$\,\% after 3--4-bit trim; $C_p\le0.85$\,pF (0.63\,pF estimated$^\dagger$); RL $\ge10$\,dB at $\fN$ & \cref{sec:termination,sec:esd-pad}\\
CTLE & peaking, poles, noise & 0--12\,dB programmable; $f_{p1}\approx2.5$\,GHz, $f_{p2}\ge7.5$\,GHz; noise $\le0.25$\,mV rms$^\dagger$ & \cref{ch:ctle}\\
DFE & taps, loop timing & half-rate, 2 taps, direct feedback; $w_1$ up to $0.5\,p_0$, $w_2$ to $\pm0.15\,p_0$; tap-1 loop $\approx$140\,ps$^\dagger$ $<$ 1\,UI; loop unrolling as fallback & \cref{sec:dfe-timing}\\
Samplers & count, sensitivity, offset & 4 HBT CML master--slave (D0 E0 D1 E1); $V_\mathrm{sens}\approx20$\,mV; offset $\sigma\approx4.9$\,mV raw$^\dagger$, $\pm0.46$\,mV after 5-bit trim; $\tau\le5$\,ps; aperture $\le0.2$\,UI & \cref{sec:sampler,sec:offset}\\
Deserializer & ratio, output & 1:16 CMOS; 16 data + 16 edge bits at 312.5\,MHz; two extra latches on the edge path; FIFO to the core & \cref{sec:deserializer,sec:metastability}\\
PLL & architecture, jitter & integer-N CP-PLL, $f_\mathrm{ref}$ = 50\,MHz, $N$ = 50; LC VCO 2.5\,GHz (4\,nH, $Q$ = 16.8); $\Kvco\approx2\pi\times150$\,MHz/V; $f_u\approx3$\,MHz; RJ $\le1.0$\,ps rms$^\dagger$ & \cref{ch:pll}\\
Quadrature & accuracy & 2-stage polyphase + CML buffers; I/Q error $\le\ang{2}$ (3$\sigma$); duty cycle $50\pm0.5$\,\% with DCC & \cref{sec:quadrature}\\
PI & resolution, linearity & 2 PIs, 64 steps per 400\,ps (6.25\,ps = 0.031\,UI); INL $\le\ang{1}$ & \cref{sec:picdr}\\
CDR & loop, tracking & half-rate Alexander BBPD; P+I at 625\,MHz, $K_P$ = 1/8 LSB, $K_P/K_I$ = 128; $\pm300$\,ppm; lock $<1$\,\textmu{}s; wander $\le20$\,ps pp$^\dagger$; JTOL above the USB-style mask & \cref{sec:bbcdr,sec:jtol}\\
Timing budget & total jitter & TJ $\le0.3$\,UI (60\,ps) at $10^{-12}$; estimated 58.6\,ps$^\dagger$ & \cref{sec:jitter-budget}\\
Voltage budget & margin at the sampler & $\ge50$\,mV after noise, offset and sensitivity; 76\,mV on REF with 2-tap de-emphasis$^\dagger$ & \cref{sec:noisebudget}\\
Test & on-chip & PRBS7/31 generator and checker; near-end loopback; SPI registers & \cref{sec:labtest}\\
Power & total, per bit & $\approx$216\,mW, 43\,pJ/b$^\dagger$ & \cref{sec:power-budget}\\
\bottomrule
\end{tabularx}
\end{table}

\paragraph{The link budget in one paragraph} Both budgets close (\cref{fig:ch14-budget}). In voltage, REF with only the 2-tap TX de-emphasis leaves a 131\,mV statistical eye at $10^{-12}$; noise ($2Q\sign=21$\,mV), jitter (3.5\,mV), offset (10\,mV) and sensitivity (20\,mV) leave a \textbf{76\,mV} margin by linear addition, 89\,mV statistically (\cref{sec:noisebudget}), and the full partition ($-6$\,dB preset + 2-tap DFE, 172\,mV noise-free PRBS7 eye, \cref{sec:eq-partition}) only adds to it. In time, 40\,ps of DJ plus $2Q\times1.32$\,ps of RJ give TJ = \textbf{58.6\,ps = 0.293\,UI}, just inside 0.3\,UI (\cref{sec:jitter-budget}). Timing is the tighter budget, and its largest line is the CDR's own dither.

\begin{figure}[!tbp]
\centering
\includegraphics[width=\linewidth]{ch14_budget}
\caption{\textbf{Both budgets close, but timing is the tight one.} (a) The 131\,mV statistical eye of REF with 2-tap de-emphasis, minus the terms of \cref{sec:noisebudget}: 76\,mV of margin remains, more than half the eye. (b) One 200\,ps UI minus the deterministic jitter (linear sum) and $2Q$ times the random jitter (RSS): 141\,ps remains, only 1.4\,ps more than the 0.7\,UI target allows.}
\label{fig:ch14-budget}
\end{figure}

%% =====================================================================
\section{Power budget}\label{sec:power-budget}

A CML gate draws its tail current whatever the data, so a CML-heavy chip has a power budget you can write down before any layout: list the tails, multiply by the supply. CMOS logic burns $\alpha C V^2 f$ instead.

\begin{keyeq}[CML power and energy per bit]
\begin{equation}
  P_\mathrm{CML} = V_\mathrm{CC}\sum_i I_{T,i},\qquad P_\mathrm{CMOS}=\alpha\,C_\mathrm{sw}V_\mathrm{DD}^2f,\qquad E_b = \frac{P_\mathrm{total}}{\fb}
  \label{eq:ch14-eb}
\end{equation}
\end{keyeq}

\Cref{tab:ch14-power} collects the estimates of \cref{ch:tx,ch:rx,ch:pll,ch:cdr} and \cref{fig:ch14-power} plots them. Assumptions: 1\,Vppd swing (20\,mA driver tail); CML gates as in \cref{we:ch9-hbt}, about 1.2\,mA each with emitter followers; CML on 2.5\,V, CMOS on 1.2\,V; the team's measured PLL power of 12.73\,mW~\cite{smacd2026pll}. The polyphase buffers, the PIs and the sampler clock buffers are separate circuits and are counted separately.

\begin{table}[!tbp]
\centering
\caption{\textbf{First-cut \OS{} power budget.} Currents from the checklists and target boxes of \cref{ch:tx,ch:rx,ch:pll,ch:cdr}; every line is an estimate to be replaced by simulation.}
\label{tab:ch14-power}
\footnotesize
\begin{tabularx}{\linewidth}{@{}l >{\raggedright\arraybackslash}p{52mm} L r@{}}
\toprule
& \textbf{Blocks} & \textbf{Assumption} & \textbf{mW}\\
\midrule
TX & 3-tap CML driver & $I_T$ = 20\,mA from 2.5\,V (40\,mW at 0.8\,Vppd) & 50.0\\
& CML 4:2, 2:1 ($\times3$), latches, pre-drivers, clock buffers, $\div2$ & about 17 CML gates of 1.2\,mA & 51.0\\
& CMOS 16:4 & 1.2\,V, 625\,MHz & 3.0\\
RX & CTLE + buffer, DFE summers and tap DACs & 5\,mA + 6\,mA & 27.5\\
& 4 CML samplers + their clock buffers & 12\,mA + 4\,mA & 40.0\\
& CMOS 1:16 deserializer, offset DACs & 1.2\,V; 5 bits $\times4$ & 5.6\\
Clocking & PLL (team design) & measured, 130\,nm BiCMOS~\cite{smacd2026pll} & 12.7\\
& polyphase + I/Q buffers, distribution, 2 PIs & 6\,mA + 8\,mW (\cref{sec:picdr}) & 23.0\\
Digital & CDR loop filter, PRBS gen.\ + checker, SPI & CMOS at 625 and 312.5\,MHz & 3.2\\
\midrule
\textbf{Total} & & TX 104, RX 73, clocking 36, digital 3 & \textbf{216}\\
\bottomrule
\end{tabularx}
\end{table}

\begin{figure}[!tbp]
\centering
\includegraphics[width=\linewidth]{ch14_power}
\caption{\textbf{The CML driver alone is almost a quarter of the chip's power; CMOS logic is a rounding error.} (a) The first-cut budget of \cref{tab:ch14-power}, colored by section. The TX takes 48\,\%, mostly static CML current; the four samplers are the largest RX item. (b) Energy per bit against the Purdue OpenSerDes, an all-digital 2\,Gb/s link in SKY130 at about 438\,mW~\cite{openserdes2021}, and against \OS{} with the voltage-mode (SST) driver of \cref{sec:cml-driver}.}
\label{fig:ch14-power}
\end{figure}

\begin{example}{Energy per bit, and where to save}\label{we:ch14-power}
\textbf{Total:} the CML tails add up to 77\,mA from 2.5\,V (192\,mW); the PLL, the CMOS blocks and the DACs add 25\,mW: $P\approx\qty{216}{\milli\watt}$. By \cref{eq:ch14-eb}, $E_b = 216\,\mathrm{mW}/5\,\mathrm{Gb/s} = \mathbf{43.2}$\,pJ/b.

\textbf{Comparison:} the Purdue OpenSerDes spends $438\,\mathrm{mW}/2\,\mathrm{Gb/s} = 219$\,pJ/b~\cite{openserdes2021}. \OS{} runs 2.5 times faster for about half the power, five times less energy per bit, because the HBTs do in a few CML gates what an all-digital 130\,nm CMOS design must do with oversampling and many more gates.

\textbf{Savings, largest first:} (1) the SST driver of \cref{sec:cml-driver} saves about 44\,mW, giving 172\,mW and 34.4\,pJ/b; (2) a swing of 0.8\,Vppd saves 10\,mW; (3) each CML gate removed from the clock paths saves about 3\,mW. Lowering the CML supply would save the most of all, but every $V_\mathrm{CE}$ must still fit below $\mathrm{BV}_\mathrm{CEO}$ with the swing on top (\cref{sec:cml-driver}), so this waits for silicon data.
\end{example}


%% =====================================================================
\section{Design decisions}\label{sec:ch14-decisions}

\Cref{tab:decisions} lists the ten choices that shape \OS{}. Each row says what was chosen, what was rejected and why, and where the argument is made in full. Use it in design reviews: if someone proposes to change a row, the chapter in the last column says which numbers must be recomputed.

\begin{table}[!tbp]
\centering
\caption{\textbf{The ten main \OS{} architecture decisions.}}
\label{tab:decisions}
\footnotesize
\renewcommand{\arraystretch}{1.12}
\begin{tabularx}{\linewidth}{@{}>{\raggedright\arraybackslash}p{17mm} >{\raggedright\arraybackslash}p{22mm} >{\raggedright\arraybackslash}p{21mm} L >{\raggedright\arraybackslash}p{16mm}@{}}
\toprule
\textbf{Decision} & \textbf{Chosen} & \textbf{Alternatives} & \textbf{Why} & \textbf{Ch.}\\
\midrule
Modulation & NRZ & PAM4 & on REF NRZ wins by 2.6\,dB; PAM4 pays only above about 24\,dB and needs 3 slicers and linear circuits & \cref{sec:pam4}\\
Clock rate & half rate, 2.5\,GHz & full rate, quarter rate & full rate needs 5\,GHz clocks; quarter rate needs 8 calibrated phases & \cref{sec:serializer,sec:quadrature}\\
Fast logic & SiGe HBT CML & CMOS, MOS CML & full switching at 100\,mV, gain 5.8 vs.\ 2; CMOS misses the 8\,FO4 rule at 2.5\,GHz, slow corner & \cref{sec:serializer}\\
TX driver & CML & SST (voltage mode) & resistor-set impedance, constant current, taps leave the impedance alone; SST ($-44$\,mW) is the upgrade & \cref{sec:cml-driver}\\
TX FFE & 3 taps & 2 taps, 4+ taps & post-tap does most of the work; pre-tap for $p_{-1}$ on longer channels (on REF, 2 taps win: Lab~3) & \cref{ch:txffe,sec:lab3}\\
RX equalizer & CTLE 0--12\,dB + 2-tap DFE & CTLE only, more DFE taps, RX FFE & DFE removes $p_1$, $p_2$ without noise; CTLE adds gain, adaptivity and reach above 22\,dB & \cref{sec:eq-partition}\\
VCO & LC, 4\,nH, $Q$ = 16.8 & ring VCO & 15--30\,dB lower phase noise; reuses the team's proven design & \cref{sec:lcvco}\\
CDR & PI dual loop, digital BB & VCO-based BB, linear (Hogge) & reuses the clean PLL; synthesizable; PVT-independent gain; no false lock & \cref{sec:picdr}\\
Quadrature & 2-stage polyphase + CML & 5\,GHz VCO + $\div2$, quadrature VCO & keeps the 2.5\,GHz VCO; $\div2$ needs a 5\,GHz VCO with an accurate duty cycle & \cref{sec:quadrature}\\
On-chip test & PRBS7/31 + checker, loopback & external BERT only & scope and PC suffice; $10^{-12}$ in 10 minutes; loopback isolates the chip & \cref{sec:labtest}\\
\bottomrule
\end{tabularx}
\end{table}

%% =====================================================================
\section{Risks and open questions}\label{sec:ch14-risks}

The architecture closes on paper. It can still fail in these places, roughly in order of the cost of a late fix:
\begin{itemize}
  \item \textbf{ESD and pad capacitance.} ESD diodes can take most of the 0.85\,pF budget alone; get the \sgx{} ESD cell capacitance early, and keep the T-coil of \cref{sec:esd-pad} as the fallback.
  \item \textbf{$\mathrm{BV}_\mathrm{CEO}$ headroom.} The off driver transistor sits at 1.25--1.5\,V against about 1.6\,V; simulate every $V_\mathrm{CE}$ at every corner with the real termination.
  \item \textbf{PLL phase noise.} The quoted $-100.8$\,dBc/Hz at 1\,MHz is far worse than Leeson predicts (\cref{sec:lcvco}); the 1\,ps RJ line depends on what it really is.
  \item \textbf{Inductor EM modeling.} The 4\,nH tank and any T-coil must be EM-simulated on TopMetal2 (\cref{sec:layout-hs}) and matched to the team's measured $Q$.
  \item \textbf{Timing loops.} DFE tap 1 uses about 140 of its 200\,ps and the CDR dither 19 of its 20\,ps; both need post-layout confirmation.
  \item \textbf{Area, pads and the channel.} The PLL is 0.347\,mm$^2$~\cite{smacd2026pll}, the MIM capacitors add tenths of a mm$^2$, and the high-speed pads need SMA-grade access; REF must be replaced by measured S-parameters (\cref{sec:ch16-next}).
\end{itemize}

\begin{takeaways}
\begin{itemize}
  \item \OS{} is a single-lane transceiver: CMOS at 312.5\,MHz in the core, HBT CML at 1.25\,GHz and above, joined by level converters, two supplies and two clock domains.
  \item One 2.5\,GHz LC PLL feeds both sides: the TX divides it; the RX makes four phases and rotates them in 6.25\,ps steps under digital CDR control.
  \item The block specification table is the contract: every block has a number and a chapter; $^\dagger$ marks estimates awaiting simulation.
  \item The voltage budget leaves about 76\,mV; the timing budget leaves a 141\,ps eye with 1.4\,ps to spare.
  \item First-cut power is about 216\,mW (43\,pJ/b), five times less energy per bit than OpenSerDes, dominated by static CML current; the SST driver is the first saving.
\end{itemize}
\end{takeaways}
